\documentclass{article}

\usepackage[a4paper, margin=2.6cm, headsep=18pt]{geometry}
\usepackage{amsmath, amssymb}
\usepackage{graphicx}
\usepackage{booktabs}
\usepackage{caption}
\usepackage[table]{xcolor}
\usepackage{titlesec}
\usepackage[hidelinks]{hyperref}
\usepackage{setspace}
\usepackage{microtype}
\usepackage{enumitem}
\usepackage{tcolorbox}

% ---------- palette ----------
\definecolor{ink}{RGB}{30,32,44}
\definecolor{indigo}{RGB}{44,54,98}
\definecolor{indigofaint}{RGB}{237,239,247}
\definecolor{indigoline}{RGB}{150,160,195}

\color{ink}

% ---------- section styling ----------
\titleformat{\section}
  {\Large\bfseries\color{indigo}}
  {\thesection}{0.6em}{}
  [{\color{indigoline}\titlerule[0.6pt]}]
\titlespacing*{\section}{0pt}{22pt}{10pt}

\titleformat{\subsection}
  {\large\bfseries\color{indigo}}
  {\thesubsection}{0.6em}{}
\titlespacing*{\subsection}{0pt}{16pt}{6pt}


% ---------- callout box for key numerical results ----------
\newtcolorbox{keyresult}{
  colback=indigofaint, colframe=indigo,
  boxrule=0.6pt, left=8pt, right=8pt, top=6pt, bottom=6pt,
  before skip=10pt, after skip=10pt
}

\setstretch{1.1}

\captionsetup{font={small}, labelfont={bf,color=indigo}, margin=10pt}

\title{\textbf{\Large Measuring Electromagnetic Induction:\\[2pt] The Peak EMF Induced by a Magnet Falling Through a Coil}}
\author{Chhaya Ramnani (u8336188)\\ \small Lab partners: Tsz Wing Candice Ng, Chinatsu Komada\\ \small PHYS1201 --- Physics 2, Week 5 Laboratory: Electromagnetic Induction}
\date{Date of submission: 3 September 2026}

\begin{document}
 
\begin{titlepage}
\centering
\vspace*{3cm}
{\color{indigoline}\rule{0.5\linewidth}{0.6pt}}\\[0.6cm]
{\Huge\bfseries\color{indigo} Measuring Electromagnetic\\[4pt] Induction}\\[10pt]
{\Large\itshape The Peak EMF Induced by a Magnet\\ Falling Through a Coil}\\[0.6cm]
{\color{indigoline}\rule{0.5\linewidth}{0.6pt}}\\[1.8cm]
{\large Chhaya Ramnani \ \textnormal{(u8336188)}}\\[6pt]
{\normalsize Lab partners: Tsz Wing Candice Ng, Chinatsu Komada}\\[1.4cm]
{\normalsize PHYS1201 --- Physics 2}\\
{\normalsize Week 5 Laboratory: Electromagnetic Induction}\\[1.2cm]
{\normalsize Date of submission: 3 September 2026}
\vfill
\end{titlepage}
 
\setcounter{page}{1}
 
\section*{1.0 Abstract}
This experiment examined how the peak electromotive force (emf) induced by a magnet falling 
through a coil depends on three things: the number of coil turns $N$, the coil radius $R$, and 
the magnet's speed $s$ as it passed through. The magnet's dipole moment was first found from calibration 
measurements of its on-axis field, giving $m = (1.18 \pm 0.07)$~A\,m$^2$. This was used to build a 
theoretical prediction for the peak emf, which the measured data from each of the three sub-experiments 
was then compared against directly: no scaling factor was fitted anywhere. The peak emf rose 
roughly linearly with $N$ and with $s$, and fell off roughly as $R^{-2}$ with coil radius, matching the 
predicted shape $\varepsilon_{\text{peak}} \propto Ns/R^2$. None of the three datasets matched 
the absolute theoretical prediction within uncertainty, and the $N$ data disagreed by far the most 
(a measured-to-theory ratio of $0.461 \pm 0.003$). This was mainly put down to the magnet's finite size 
breaking the point-dipole assumption the theory relies on, and, specifically for the $N$ sub-experiment, 
to an unresolved mismatch between two setups that should have given the same reading.
 
\section*{2.0 Introduction}
 
\subsection*{2.1 Aims}
This experiment set out to see how the peak emf induced by a magnet falling through a coil depends on three 
parameters: the number of turns in the coil, $N$; the coil radius, $R$; and the magnet's speed as it passed 
through, $s$. A dipole model of the magnet was calibrated first from independent field measurements, and the 
resulting theoretical prediction was then compared directly against the measured emf for each sub-experiment.
 


\end{document}